\section{When to Use Each Variant}
\label{sec:when-to-use}

The three \sburst{} variants serve different purposes.
The choice depends on whether the primary goal is optimization (find the
best plan) or characterisation (describe the distribution of plans near
the optimum).

\subsection{Standard Short-Burst: Pure Compactness Optimization}

Use \textbf{SB-Standard} (\texttt{--search short-burst}) when:
\begin{itemize}
  \item The goal is to find a single plan with minimum edge-cut.
  \item Compute budget is the binding constraint.
  \item Distributional correctness of the endpoint is not required.
\end{itemize}

At equal compute budget, SB-Standard visits the most distinct plans per
burst (highest $\bar\alpha$) and consistently achieves the lowest final
edge-cut.
It is the fastest variant and the best choice for production plan generation.

\paragraph{Typical parameters.}
\begin{verbatim}
redist state --state TX --year 2020 \
    --search short-burst \
    --burst-length 20 --n-bursts 100 --percentile 0.0
\end{verbatim}

\subsection{ShortBurstForest: Calibrated Exploration Near Good Plans}

Use \textbf{\sbf{}} (\texttt{--search short-burst-forest}) when:
\begin{itemize}
  \item The goal is to characterise the distribution of partisan outcomes
        for plans with edge-cut within a specified range of the optimum.
  \item The reported ensemble will be used in legislative findings, expert
        reports, or litigation, where the claim is that ``among all plans
        within $X$\% of the optimal edge-cut, $Y$\% produce outcome $O$''.
  \item Approximate distributional correctness is required to support
        statistical claims about the ensemble.
\end{itemize}

\sbf{} produces endpoints that are approximately exchangeable draws from
the stationary distribution of \frecom{} within the burst neighbourhood of
the current plan.
This is the appropriate choice when the distribution of the endpoints
matters, not just the minimum.

\paragraph{Typical parameters.}
\begin{verbatim}
redist state --state NC --year 2020 \
    --search short-burst-forest \
    --burst-length 20 --n-bursts 50 --percentile 0.0
\end{verbatim}

\subsection{ShortBurstMergeSplit: The Practical Middle}

Use \textbf{\sbms{}} (\texttt{--search short-burst-merge-split}) when:
\begin{itemize}
  \item Distributional correctness is desired but \sbf{}'s compute cost
        (two RNG streams per step) is prohibitive.
  \item A moderate improvement in distributional fidelity relative to
        SB-Standard is acceptable at modest cost in final EC.
  \item The state has large $k$ (e.g., $k \geq 30$), where the acceptance-rate
        gap between \sbms{} and SB-Standard is smallest.
\end{itemize}

\sbms{} offers intermediate acceptance rate ($\bar\alpha \approx 0.60$)
and intermediate EC quality.
It is cheaper than \sbf{} (single MH ratio per step) while providing
more distributional correction than SB-Standard (no correction).

\paragraph{Typical parameters.}
\begin{verbatim}
redist state --state TX --year 2020 \
    --search short-burst-merge-split \
    --burst-length 20 --n-bursts 50 --percentile 0.0
\end{verbatim}

\subsection{Audit Trail}

All three variants record the following fields in the run audit manifest:
\begin{itemize}
  \item \texttt{burst\_seeds}: the list of $s_i$ values, one per burst.
  \item \texttt{selected\_burst\_idx}: the index $i^*$ of the selected
        endpoint (the burst with the minimum EC at percentile $p$).
  \item \texttt{per\_burst\_acceptance\_rates}: the list of $\alpha_i$
        values, one per burst.
  \item \texttt{chain\_type}: the string identifier of the inner chain
        (\texttt{"recom"}, \texttt{"forest-recom"}, or
        \texttt{"merge-split"}).
\end{itemize}

An independent verifier can reproduce the trajectory of any burst by:
\begin{enumerate}
  \item Reconstructing $s_i$ from the base seed $b$, burst index $i$,
        and chain-type prefix.
  \item Re-running the inner chain for $\ell$ steps using the per-step
        forward and reverse seeds derived from $s_i$.
  \item Verifying that the endpoint assignment matches the recorded plan.
\end{enumerate}

The \texttt{per\_burst\_acceptance\_rates} field also serves as an
integrity check: if the verifier's replication produces different
$\alpha_i$ values, the seed derivation or chain implementation is
inconsistent.

\subsection{Compute Budget Guidance}

The equal-budget comparison ($\ell = 20$, $n = 50$, total $= 1000$ steps)
is a baseline.
Recommended adjustments by state size:
\begin{itemize}
  \item Small states ($k \leq 10$): $\ell = 10$, $n = 100$ (more bursts
        improve diversity at fixed total).
  \item Medium states ($10 < k \leq 20$): $\ell = 20$, $n = 50$ (baseline).
  \item Large states ($k > 20$): $\ell = 30$, $n = 50$--$100$ (longer
        bursts needed to explore meaningfully within small districts).
\end{itemize}

For \sbf{} specifically, the effective burst length is
$\bar\alpha \cdot \ell \approx 0.45\,\ell$.
To match SB-Standard's effective per-burst exploration, use
$\ell \approx 33$ with \sbf{}---a 65\% increase in nominal steps.
Practitioners should account for this when budgeting compute time.
